% GENERATED FILE. Edit canonical lessons, then run npm run generate:books.
% canonical-source-hash: fnv1a64:5ed49ddd733f6865
% canonical-lessons: FR-C141-lavenir, FR-C141-ce-soir, FR-C141-le-projet, FR-C141-prevoir, FR-C141-organiser

\chapter{The Future, This Evening, a Plan, To Foresee, To Organise}
\label{ch:fr-the-future-this-evening-a-plan-to-foresee-to-organise}
\hlchaptermodality{141}

\begin{chapteropening}
\textbf{By the end of this chapter:} \emph{I can say what will happen, and when: this evening, next week, in a year.}

Complete the last lesson of chapter 141: I can say what will happen, and when: this evening, next week, in a year.
\end{chapteropening}

\section[l'avenir]{l'avenir — the future}

\label{lesson:FR-C141-lavenir}



\begin{quote}
Before the new one: say the French for formerly, then the French for to happen.
\end{quote}

\subsection*{You'll want to know: l'avenir}
\textbf{l'avenir} — \textquotedblleft{}the future\textquotedblright{}. French builds the future on the infinitive: \textbf{parler} becomes \textbf{je parlerai} — \textquotedblleft{}I will speak\textquotedblright{} — then \textbf{tu parleras}, \textbf{il parlera}.

\begin{grammarlens}[title={What you've built}]
One more word for this chapter.
\end{grammarlens}

\subsection*{Your turn}
\begin{itemize}
\raggedright
  \item \emph{Say it:} \emph{l'avenir}
  \item \emph{Say it:} \emph{l'avenir}, once more
  \item \emph{Say it:} say \emph{l'avenir}
  \item \emph{Recall:} say the French for formerly, then the French for to happen, then say \emph{l'avenir} again
\end{itemize}

\begin{tcolorbox}[breakable,colback=teal!4,colframe=teal!35!black,title={Before you move on}]
What does l'avenir mean? (\textquotedblleft{}The future\textquotedblright{}.) Say it once more.
\end{tcolorbox}
\section[ce soir]{ce soir — this evening}

\label{lesson:FR-C141-ce-soir}



\begin{quote}
Before the new one: say the French for to happen, then the French for the future.
\end{quote}

\subsection*{You'll want to know: ce soir}
\textbf{ce soir} — \textquotedblleft{}this evening\textquotedblright{}. The time word can stand first: \textbf{Ce soir, je travaillerai.} — \textquotedblleft{}This evening I will work.\textquotedblright{} \textbf{Dans} a length of time gives \textquotedblleft{}in\textquotedblright{}: \textbf{dans deux jours} — \textquotedblleft{}in two days\textquotedblright{}.

\begin{grammarlens}[title={What you've built}]
One more word for this chapter.
\end{grammarlens}

\subsection*{Your turn}
\begin{itemize}
\raggedright
  \item \emph{Say it:} \emph{ce soir}
  \item \emph{Say it:} \emph{ce soir}, once more
  \item \emph{Say it:} say \emph{ce soir}
  \item \emph{Recall:} say the French for to happen, then the French for the future, then say \emph{ce soir} again
\end{itemize}

\begin{tcolorbox}[breakable,colback=teal!4,colframe=teal!35!black,title={Before you move on}]
What does ce soir mean? (\textquotedblleft{}This evening\textquotedblright{}.) Say it once more.
\end{tcolorbox}
\section[le projet]{le projet — a plan, a project}

\label{lesson:FR-C141-le-projet}



\begin{quote}
Before the new one: say the French for the future, then the French for this evening.
\end{quote}

\subsection*{You'll want to know: le projet}
\textbf{le projet} — \textquotedblleft{}a plan, a project\textquotedblright{}. \textbf{Quels sont tes projets pour le week-end ?} — \textquotedblleft{}What are your plans for the weekend?\textquotedblright{}

\begin{grammarlens}[title={What you've built}]
One more word for this chapter.
\end{grammarlens}

\subsection*{Your turn}
\begin{itemize}
\raggedright
  \item \emph{Say it:} \emph{le projet}
  \item \emph{Say it:} \emph{le projet}, once more
  \item \emph{Say it:} say \emph{le projet}
  \item \emph{Recall:} say the French for the future, then the French for this evening, then say \emph{le projet} again
\end{itemize}

\begin{tcolorbox}[breakable,colback=teal!4,colframe=teal!35!black,title={Before you move on}]
What does le projet mean? (\textquotedblleft{}A plan, a project\textquotedblright{}.) Say it once more.
\end{tcolorbox}
\section[prévoir]{prévoir — to foresee, to plan for}

\label{lesson:FR-C141-prevoir}



\begin{quote}
Before the new one: say the French for this evening, then the French for a plan.
\end{quote}

\subsection*{You'll want to know: prévoir}
\textbf{prévoir} — \textquotedblleft{}to foresee, to plan for\textquotedblright{}. \textbf{voir} with \textbf{pré-}, \textquotedblleft{}before\textquotedblright{}: seeing a thing before it comes. The future of \textbf{être} is irregular: \textbf{je serai}, \textbf{tu seras}, \textbf{il sera}.

\begin{grammarlens}[title={What you've built}]
One more word for this chapter.
\end{grammarlens}

\subsection*{Your turn}
\begin{itemize}
\raggedright
  \item \emph{Say it:} \emph{prévoir}
  \item \emph{Say it:} \emph{prévoir}, once more
  \item \emph{Say it:} say \emph{prévoir}
  \item \emph{Recall:} say the French for this evening, then the French for a plan, then say \emph{prévoir} again
\end{itemize}

\begin{tcolorbox}[breakable,colback=teal!4,colframe=teal!35!black,title={Before you move on}]
What does prévoir mean? (\textquotedblleft{}To foresee, to plan for\textquotedblright{}.) Say it once more.
\end{tcolorbox}
\section[organiser]{organiser — to organise}

\label{lesson:FR-C141-organiser}



\begin{quote}
Before the new one: say the French for a plan, then the French for to foresee.
\end{quote}

\subsection*{You'll want to know: organiser}
\textbf{organiser} — \textquotedblleft{}to organise\textquotedblright{}. A regular \textbf{-er} verb, so the future follows \textbf{parler}: \textbf{J'organiserai tout.} — \textquotedblleft{}I will organise everything.\textquotedblright{}

\begin{grammarlens}[title={What you've built}]
One more word for this chapter.
\end{grammarlens}

\subsection*{Your turn}
\begin{itemize}
\raggedright
  \item \emph{Say it:} \emph{organiser}
  \item \emph{Say it:} \emph{organiser}, once more
  \item \emph{Say it:} say \emph{organiser}
  \item \emph{Recall:} say the French for a plan, then the French for to foresee, then say \emph{organiser} again
\end{itemize}

\begin{tcolorbox}[breakable,colback=teal!4,colframe=teal!35!black,title={Before you move on}]
What does organiser mean? (\textquotedblleft{}To organise\textquotedblright{}.) Say it once more.
\end{tcolorbox}
